\scanpage{12}\mspage{11}
Les éléments irréductibles de $\unsure{\mathcal F|\mathscr F}$ sont les $\Phi\in\mathcal F$ qui ont \deleted{\unsure{une}} un plus grand élément : car un \unsure{ens} fermé est irréductible (au sens fort) \unsure{sss|ssi} il a un plus grand élément.

% slanted
\margin{Appelons-les les ``\unsure{multistrates} ensemblistes''}

Donc

\uline{C 2} Tout $\Phi\in\deleted{\unsure{\Phi}}\mathcal F$ est la sup des objets irréductibles de $\mathcal F$, à savoir des $\Phi_{X}=\{Y\in\Phi\mid Y\unsure{\subseteq|\leq}X\}$, pour $Y\in\Phi$

\uline{C 3} Soient $\deleted{\Phi}F,\deleted{\Psi}G\in\mathcal F$. Pour que $F$ et $G$ soient compatibles (i.e $F,G$ \unsure{majoré|majorés}) il f. et s. que $\deleted{X\in F,\unsure{\Rightarrow|\to}}$ $\forall\,\deleted{\Phi}\unsure{X|{\Phi}'}\leq F$ et $\deleted{\Psi}Y\leq G$, $\deleted{\Phi}X$ et $\deleted{\Psi}Y$ \deleted{\unsure{élé|él.}} multistrates, on ait $\deleted{\Phi}\deleted{\Phi}X$ et $\deleted{\Psi}Y$ compatibles.

En d'autres termes, que pour $X\in F,\ Y\in G$, $F_X$ et $G_Y$ soient compatibles :

\uline{\unsure{Remords|Remonter|Remontons} à prop. page \deleted{\unsure{9}}7} : Cond. équiv à a) b) c)

d) \deleted{\unsure{lorsque|toujours|comp.}} $\forall\,A\in\Phi,\ B\in\Psi$, $A\cap B$ est fermé dans $A$ (pour $\Phi$ i.e pour $\Phi_A$) et en $B$ (pour $\Psi$ i.e. pour $\Psi_B$), et si $\deleted{\unsure{A}\ C\in\Phi\cup\Psi}$ $C$ est une partie de $A\cap B$, alors $C\in\Phi$ ($\Leftrightarrow C\in\Phi_A$) sss $C\in\Psi$ ($\Leftrightarrow C\in\Psi_B$) i.e $\Phi_{A\cap B}=\Psi_{A\cap B}$
